\documentclass[11pt,twoside,openany]{memoir}
\usepackage[T1]{fontenc}
\usepackage[hidelinks]{hyperref}
\usepackage{amsmath}
\usepackage[fixamsmath]{mathtools}
\usepackage{amsthm}
\usepackage{amssymb}
\usepackage[margin=1in]{geometry}
\usepackage{fancyhdr}
    \fancyhf{}

    \pagestyle{fancy}
    \cfoot{\scriptsize \thepage}
    \fancyhead[L]{\scalebox{0.8}{Algebra Test Bank}}
    \fancyhead[R]{\scalebox{0.8}{\leftmark, \rightmark}}

    \fancypagestyle{plain}{
        \fancyhf{}
        \cfoot{\scriptsize \thepage}
        \fancyhead[L]{\scalebox{0.8}{Algebra Test Bank}}
        \fancyhead[R]{\scalebox{0.8}{\leftmark, \rightmark}}
    }
\usepackage{thmtools}
    \declaretheoremstyle[
        spaceabove=10pt,
        spacebelow=10pt,
        headfont=\normalfont\bfseries,
        notefont=\mdseries, notebraces={(}{)},
        bodyfont=\normalfont,
        postheadspace=0.5em
    ]{defs}

    \declaretheoremstyle[ 
        spaceabove=10pt, 
        spacebelow=10pt,
        headfont=\normalfont\bfseries,
        bodyfont=\normalfont\itshape,
        postheadspace=0.5em
    ]{thmstyle}
    
    \declaretheorem[
        style=thmstyle,
    ]{theorem}

    \declaretheorem[
        style=thmstyle,
        sibling=theorem,
    ]{proposition}

    \declaretheorem[
        style=thmstyle,
    ]{lemma}

    \declaretheorem[
        style=thmstyle,
        sibling=theorem,
    ]{corollary}

    \declaretheorem[
        style=defs,
    ]{example}

    \declaretheorem[
        style=defs,
    ]{definition}

    \declaretheoremstyle[
        spaceabove=15pt,   
        spacebelow=15pt,   
        headfont=\normalfont\bfseries,
        notefont=\mdseries, notebraces={(}{)},
        bodyfont=\normalfont,
        postheadspace=0.5em
    ]{exerciseStyle}

    \declaretheorem[
        style=exerciseStyle,
        sibling=theorem
    ]{exercise}

    \declaretheorem[
        numbered=unless unique,
        shaded={rulecolor=black,
        rulewidth=1pt, bgcolor={rgb}{1,1,1}}
    ]{axiom}

    \declaretheorem[numberwithin=section,style=defs]{note}
    \declaretheorem[numbered=no,style=defs]{question}
    \declaretheorem[numbered=no,style=defs]{recall}
    \declaretheorem[numbered=no,style=remark]{answer}
    \declaretheorem[numbered=no,style=remark]{solution}
    \declaretheorem[numbered=no,style=defs]{remark}
\usepackage{enumitem}
\usepackage{titlesec}
    \titleformat{\chapter}[display]
    {\bfseries\Huge\raggedright}
    {Chapter {\thechapter}}
    {1ex minus .1ex}
    {\Huge}
    \titlespacing{\chapter}
    {3pc}{*3}{15pt}[3pc]

    \titleformat{\section}[block]
    {\normalfont\bfseries\LARGE}
    {\S\ \thesection.}{.5em}{}[]
    \titlespacing{\section}
    {0pt}{3ex plus .1ex minus .2ex}{3ex plus .1ex minus .2ex}
\usepackage[utf8x]{inputenc}
\usepackage{tikz}
\usepackage{tikz-cd}
\usepackage{wasysym}
\linespread{1.00}
%%%%%%%%%%%%%%%%%%%%%%%%%%%%%%%%%%%%%%%%%%%%%%%%%%%%%%%%%%%%%
%%%%%%%%%%%%%%%%%%%%%%%%%%%%%%%%%%%%%%%%%%%%%%%%%%%%%%%%%%%%%
\begin{document}

\begin{abstract}
    This bank contains two types of problems: template problems and pool problems. \textbf{Template problems} are generally computational with easily adjustable specifics. These types of problems are especially prevalent in linear algebra. \textbf{Pool problems} include all problems that are not easily adjustable. These problems, when chosen, will usually be asked as is.
\end{abstract}

\begingroup
\let\clearpage\relax
\chapter*{Linear Algebra}
\section*{\S\ Pool Problems}
\markboth{Linear Algebra}{Pool Problems}
\endgroup
    \begin{exercise}
        \phantom{a}
        \begin{enumerate}[label = (\arabic*)]
            \item Give an explicit example (with proof) showing that the union of two subspaces (of a given vector space) is not necessarily a subspace.
            \item Suppose $U_1$ and $U_2$ are subspaces of a vector space $V$. Recall that their \textbf{sum} is defined to be the set $U_1+U_2 = \left\{ u_1 + u_2 \mid u_1 \in U_1, u_2 \in U_2 \right\}$. Prove $U_1 + U_2$ is a subspace of $V$ containing $U_1$ and $U_2$.
        \end{enumerate}
        \end{exercise}
        
        \begin{exercise}
        Suppose $F$ is a field and $A$ is an $n\times n$ matrix over $F$. Suppose further that $A$ possesses distinct eigenvalues $\lambda_1$ and $\lambda_2$ with $\dim \operatorname{Null}(A-\lambda_1 I_n)=n-1$. Prove $A$ is diagonalizable.
        \end{exercise}
        
        \begin{exercise}
        Let $\phi: V \to W$ be a surjective linear transformation of finite-dimensional linear spaces. Show that there is a $U \subseteq V$ such that $V = \ker(\phi) \oplus U$ and $\phi\mid_U : U \to W$ is an isomorphism. (Note that $V$ is not assumed to be an inner-product space; also note that $\ker(\phi)$ is sometimes referred to as the \textbf{null space} of $\phi$; finally, $\phi\mid_U$ denotes the restriction of $\phi$ to $U$.)
        \end{exercise}
        
        \begin{exercise}
        Suppose $V$ is a finite-dimensional real vector space and $T: V \to V$ is a linear transformation. Prove that $T$ has at most $\dim(\operatorname{range} \, T)$ distinct nonzero eigenvalues.
        \end{exercise}
        
        \begin{exercise}
        Let $T: V \to V$ be a linear transformation on a finite-dimensional vector space. Prove that if $T^2 = T$, then
        \[
        V = \ker(T) \oplus \operatorname{im}(T).
        \]
        \end{exercise}
        
        \begin{exercise}
        Let $\mathbf{R}^3$ denote the $3$-dimensional vector space, and let $\mathbf{v}=(a,b,c)$ be a fixed nonzero vector. The maps $C: \mathbf{R}^3 \to \mathbf{R}^3$ and $D: \mathbf{R}^3 \to \mathbf{R}^3$ defined by $C(\mathbf{w}) = \mathbf{v} \times \mathbf{w}$ and $D(\mathbf{w}) = (\mathbf{v} \cdot \mathbf{w}) \mathbf{v}$ are linear transformations.
        \begin{enumerate}[label = (\arabic*)]
            \item Determine the eigenvalues of $C$ and $D$.
            \item Determine the eigenspaces of $C$ and $D$ as subspaces of $\mathbf{R}^3$, in terms of $a, b, c$.
            \item Find a matrix for $C$ with respect to the standard basis.
        \end{enumerate}
        Show all work and explain reasoning.
        \end{exercise}
        
        \begin{exercise}
        Suppose $A$ is a real $n \times n$ matrix that satisfies $A^2 \mathbf{v} = 2 A \mathbf{v}$ for every $\mathbf{v} \in \mathbf{R}^n$.
        \begin{enumerate}[label = (\arabic*)]
            \item Show that the only possible eigenvalues of $A$ are 0 and 2.
            \item For each $\lambda \in \mathbf{R}$, let $E_\lambda$ denote the $\lambda$-eigenspace of $A$, i.e., $E_\lambda = \{ \mathbf{v} \in \mathbf{R}^n \mid A \mathbf{v} = \lambda \mathbf{v} \}$. Prove that $\mathbf{R}^n = E_0 \oplus E_2$. \textit{Hint:} For every vector $\mathbf{v}$ one can write $\mathbf{v} = (\mathbf{v} - \frac{1}{2} A \mathbf{v}) + \frac{1}{2} A \mathbf{v}$.
        \end{enumerate}
        \end{exercise}
        
        \begin{exercise}
        Suppose $T: \mathbf{R}^n \to \mathbf{R}^n$ is a linear transformation with distinct eigenvalues $\lambda_1, \lambda_2, \ldots, \lambda_m$, and let $\mathbf{v}_1, \mathbf{v}_2, \ldots, \mathbf{v}_m$ be corresponding eigenvectors. Prove $\mathbf{v}_1, \mathbf{v}_2, \ldots, \mathbf{v}_m$ are linearly independent.
        \end{exercise}
        
        \begin{exercise}
        Let $S: V \to V$ and $T: V \to V$ be linear transformations that commute, i.e., $S \circ T = T \circ S$. Let $v \in V$ be an eigenvector of $S$ such that $T(v) \neq 0$. Prove that $T(v)$ is also an eigenvector of $S$.
        \end{exercise}
        
        \begin{exercise}
        Suppose $A$ is a $5 \times 5$ matrix and $v_1, v_2, v_3$ are eigenvectors of $A$ with distinct eigenvalues. Prove $\{v_1, v_2, v_3\}$ is a linearly independent set. (\textit{Hint:} Consider a minimal linear dependence relation.)
        \end{exercise}
        
        \begin{exercise}
        Suppose $V$ is a vector space, and $\mathbf{v}_1, \mathbf{v}_2, \ldots, \mathbf{v}_n$ are in $V$. Prove that either $\mathbf{v}_1, \ldots, \mathbf{v}_n$ are linearly independent, or there exists a number $k \leq n$ such that $\mathbf{v}_k$ is a linear combination of $\mathbf{v}_1, \ldots, \mathbf{v}_{k-1}$.
        \end{exercise}
        
        \begin{exercise}
        Let $M_4(\mathbf{R})$ denote the 16-dimensional real vector space of $4 \times 4$ matrices with real entries, in which the vectors are represented as matrices. Let $T: M_4(\mathbf{R}) \to M_4(\mathbf{R})$ be the linear transformation defined by $T(A) = A - A^\top$.
        \begin{enumerate}[label = (\arabic*)]
            \item Determine the dimension of $\operatorname{ker}(T)$.
            \item Determine the dimension of $\operatorname{im}(T)$.
        \end{enumerate}
        \end{exercise}
        
        \begin{exercise}
        Let $A$ be a real $n \times n$ matrix and let $A^\top$ denote its transpose.
        \begin{enumerate}[label = (\arabic*)]
            \item Prove that $(A \mathbf{v}) \cdot \mathbf{w} = \mathbf{v} \cdot (A^\top \mathbf{w})$ for all vectors $\mathbf{v}, \mathbf{w} \in \mathbf{R}^n$. \textit{Hint:} Recall that the dot product $\mathbf{u} \cdot \mathbf{v}$ equals the matrix product $\mathbf{u}^\top \mathbf{v}$.
            \item Suppose now $A$ is also symmetric, i.e., that $A^\top = A$. Also suppose $\mathbf{v}$ and $\mathbf{w}$ are eigenvectors of $A$ with different eigenvalues. Prove that $\mathbf{v}$ and $\mathbf{w}$ are orthogonal.
        \end{enumerate}
        \end{exercise}
        
        \begin{exercise}
        A real $n \times n$ matrix $A$ is called \textbf{skew-symmetric} if $A^\top = -A$. Let $V_n$ be the set of all skew-symmetric matrices in $\operatorname{M}_n(\mathbf{R})$. Recall that $\operatorname{M}_n(\mathbf{R})$ is an $n^2$-dimensional $\mathbf{R}$-vector space with standard basis $\{ e_{ij} \mid 1 \le i,j \le n\}$, where $e_{ij}$ is the $n \times n$ matrix with a 1 in the $(i,j)$-position and zeros everywhere else.
        \begin{enumerate}[label = (\arabic*)]
            \item Show $V_n$ is a subspace of $\operatorname{M}_n(\mathbf{R})$.
            \item Find an ordered basis $\mathcal{B}$ for the space $V_3$ of all skew-symmetric $3 \times 3$ matrices.
        \end{enumerate}
        \end{exercise}

%%%%%%%%%%%%%%%%%%%%%%%%%%%%%%%%%%%
%%%%%%%%%%%%%%%%%%%%%%%%%%%%%%%%%%%
%%%%%%%%%%%%%%%%%%%%%%%%%%%%%%%%%%%
%%%%%%%%%%%%%%%%%%%%%%%%%%%%%%%%%%%
%%%%%%%%%%%%%%%%%%%%%%%%%%%%%%%%%%%
%%%%%%%%%%%%%%%%%%%%%%%%%%%%%%%%%%%
\newpage
\section*{\S\ Template Problems}
\markboth{Linear Algebra}{Template Problems}
\begin{exercise}
    Let $T:\mathbf{R}^3\to \mathbf{R}^3$ be the linear transformation that rotates counterclockwise around the $z$-axis by $\frac{2\pi}{3}$.
    \begin{enumerate}[label = (\arabic*)]
        \item Write the matrix for $T$ with respect to the standard basis $\left\{\begin{bmatrix} 1 \\ 0 \\ 0 \end{bmatrix},\begin{bmatrix} 0 \\ 1 \\ 0\end{bmatrix},\begin{bmatrix} 0 \\ 0 \\ 1\end{bmatrix}\right\}$.
        \item Write the matrix for $T$ with respect to the basis $\left\{\begin{bmatrix} \frac{\sqrt{3}}{2} \\ -\frac{1}{2} \\ 0 \end{bmatrix},\begin{bmatrix} 0 \\ 1 \\ 0\end{bmatrix},\begin{bmatrix} 0 \\ 0 \\ 1\end{bmatrix}\right\}$.
        \item Determine all (complex) eigenvalues of $T$.
        \item Is $T$ diagonalizable over $\mathbf{C}$? Justify your answer.
    \end{enumerate}
    \end{exercise}
    
    \begin{exercise}
    Let $V$ denote the real vector space of polynomials in $x$ of degree at most 3. Let $\mathcal{B}=\{1, x, x^2, x^3\}$ be a basis for $V$ and $T:V\to V$ be the function defined by $T(f(x))=f(x)+f'(x)$.
    \begin{enumerate}[label = (\arabic*)]
        \item Prove that $T$ is a linear transformation.
        \item Find $[T]_{\mathcal{B}}$, the matrix representation for $T$ in terms of the basis $\mathcal{B}$.
        \item Is $T$ diagonalizable? If yes, find a matrix $A$ so that $A[T]_{\mathcal{B}}A^{-1}$ is diagonal, otherwise explain why $T$ is not diagonalizable.
    \end{enumerate}
    \end{exercise}
    
    \begin{exercise}
    Let $\operatorname{M}_n(\mathbf{R})$ be the vector space of all $n \times n$ matrices with real entries. We say that $A, B \in \operatorname{M}_n(\mathbf{R})$ commute if $AB = BA$.
    \begin{enumerate}[label = (\arabic*)]
        \item Fix $A \in \operatorname{M}_n(\mathbf{R})$. Prove that the set of all matrices in $\operatorname{M}_n(\mathbf{R})$ that commute with $A$ is a subspace of $\operatorname{M}_n(\mathbf{R})$.
        \item Let $A=\begin{bmatrix} 1 & 1 \\ 1 & 1  \end{bmatrix}\in \operatorname{M}_2(\mathbf{R})$ and let $W\subseteq \operatorname{M}_2(\mathbf{R})$ be the subspace of all matrices that commute with $A$. Find a basis of $W$. 
    \end{enumerate}
    \end{exercise}
    
    \begin{exercise}
    Let $V\subset \mathbf{R}^5$ be the subspace defined by the equation
    \[
    x_1-2x_2+3x_3-4x_4+5x_5=0.
    \]
    \begin{enumerate}[label = (\arabic*)]
        \item Find (with justification) a basis for $V$.
        \item Find (with justification) a basis for $V^{\perp}$, the subspace of $\mathbf{R}^5$ orthogonal to $V$ under the usual dot product.
    \end{enumerate}
    \end{exercise}
    
    \newpage
    \begin{exercise}
    Let $T:\mathbf{R}^3\to \mathbf{R}^3$ be the linear transformation defined by
    \[
    T\left(\begin{bmatrix} x \\ y \\ z\end{bmatrix}\right) = \begin{bmatrix} x+y \\ 2z-x \\ y+2z\end{bmatrix}.
    \]
    \begin{enumerate}[label = (\arabic*)]
        \item Find the matrix that represents $T$ with respect to the standard basis for $\mathbf{R}^3$.
        \item Find a basis for the kernel of $T$.
        \item Determine the rank of $T$.
    \end{enumerate}
    \end{exercise}
    
    \begin{exercise}
    Let $A=\begin{bmatrix} 0 & 0 & -2 \\ 1 & 2 & 1 \\ 1 & 0 & 3\end{bmatrix}$.
    \begin{enumerate}[label = (\arabic*)]
        \item Determine whether $A$ is diagonalizable, and if so, give its diagonal form along with a diagonalizing matrix.
        \item Compute $A^{42}$. Remember to show all work.
    \end{enumerate}
    \end{exercise}
    
    \begin{exercise}
    Let $A=\begin{bmatrix} 2 & -1 & -1 \\ 1 & 0 & -1 \\ 1 & -1 & 0\end{bmatrix}$.
    \begin{enumerate}[label = (\arabic*)]
        \item Compute the characteristic polynomial $p_A(x)$ of $A$. It has integer roots.
        \item For each eigenvalue $\lambda$ of $A$, find a basis for the eigenspace $E_{\lambda}$.
        \item Determine if $A$ is diagonalizable. If so, give matrices $P$ and $B$ such that $P^{-1}AP=B$ and $B$ is diagonal. If not, explain carefully why $A$ is not diagonalizable.
    \end{enumerate}
    \end{exercise}
    
    \begin{exercise}
    Let $A=\begin{bmatrix} 6 & -2 & -1 \\ 10 & -3 & -2 \\ 0 & 0 & 1\end{bmatrix}$.
    \begin{enumerate}[label = (\arabic*)]
        \item Find bases for the eigenspaces of $A$.
        \item Determine if $A$ is diagonalizable. If so, give an invertible matrix $P$ and diagonal matrix $D$ such that $P^{-1}AP=D$. If not, explain why not.
    \end{enumerate}
    \end{exercise}
    
    \begin{exercise}
    Let $W\subset \mathbf{R}^5$ be the subspace spanned by the set of vectors $$\{\langle 1,-2,0,2,-1\rangle,\langle -2,4,-1,1,2\rangle,\langle 0,1,2,-2,1\rangle\}.$$
    \begin{enumerate}[label = (\arabic*)]
        \item Compute the dimension of $W$.
        \item Determine the dimension of $W^\perp$, the perpendicular subspace in $\mathbf{R}^5$.
        \item Find a basis for $W^\perp$.
    \end{enumerate}
    \end{exercise}
    
    \newpage
    \begin{exercise}
    Let $P_3$ be the real vector space of all real polynomials of degree three or less. Define $L:P_3\to P_3$ by $L(p(x))=p(x)+p(-x)$.
    \begin{enumerate}[label = (\arabic*)]
        \item Prove $L$ is a linear transformation.
        \item Find a basis for the null space of $L$.
        \item Compute the dimension of the image of $L$.
    \end{enumerate}
    \end{exercise}
    
    \begin{exercise}
    Let $V=\{a_0+a_1\sqrt[3]{2}+a_2\sqrt[3]{4}\mid a_0, a_1, a_2\in \mathbf{Q}\}\subseteq \mathbf{R}$. This set is a vector space over $\mathbf{Q}$.
    \begin{enumerate}[label = (\arabic*)]
        \item Verify $V$ is closed under product (using the usual product operation in $\mathbf{R}$).
        \item Let $T:V\to V$ be the linear transformation defined by $T(v)=(\sqrt[3]{2}+\sqrt[3]{4}) v$. Find the matrix that represents $T$ with respect to the basis $\{1,\sqrt[3]{2},\sqrt[3]{4}\}$ for $V$.
        \item Determine the characteristic polynomial for $T$.
    \end{enumerate}
    \end{exercise}

    \begin{exercise}
        Suppose $\{\mathbf{v}_1,\mathbf{v}_2,\mathbf{v}_3\}$ is a basis for $\mathbf{R}^3$ and $T:\mathbf{R}^3\to \mathbf{R}^3$ is a linear transformation satisfying the following:
            \begin{equation*}
            \begin{split}
                T(\mathbf{v}_1) &= 4\mathbf{v}_1+2\mathbf{v}_2 \\
                T(\mathbf{v}_2) &= 5\mathbf{v}_2 \\
                T(\mathbf{v}_3) = -2\mathbf{v}_1+4\mathbf{v}_2+5\mathbf{v}_3.
            \end{split}
            \end{equation*}
        Determine the eigenvalues of $T$ and find a basis for each eigenspace.
        \end{exercise}
        
        \begin{exercise}
        Let $W\subset \mathbf{R}^5$ be the space spanned by the vectors
        \[
        \left\{\begin{bmatrix} 1 \\ -2 \\ 0 \\ 2 \\ 1 \end{bmatrix},\begin{bmatrix} -2 \\ 4 \\ -1 \\ 1 \\ 2\end{bmatrix}, \begin{bmatrix} 0 \\ 1 \\ 2 \\ -2 \\1\end{bmatrix}\right\}.\]

        \begin{enumerate}[label = (\arabic*)]
            \item Compute the dimension of $W$.
            \item Let $W^{\perp}=\{\mathbf{v}\in \mathbf{R}^5\mid \mathbf{v}\cdot \mathbf{w}=0 \text{ for all } \mathbf{w}\in W\}$. Determine the dimension of $W^{\perp}$, and explain why this follows immediately from (1) using a theorem.
            \item Find a basis for $W^{\perp}$.
        \end{enumerate}
        \end{exercise}
        
        \begin{exercise}
        Let $L$ be the line in $\mathbf{R}^2$ defined by $y=-3x$, and let $T:\mathbf{R}^2\to \mathbf{R}^2$ be the linear transformation that orthogonally projects onto $L$ and then stretches along $L$ by a factor of two.
        \begin{enumerate}[label = (\arabic*)]
            \item Find the eigenvalues and an eigenbasis $\mathcal{B}$ for $T$.
            \item Determine the matrix for $T$ with respect to the basis $\mathcal{B}$.
            \item Determine the matrix for $T$ with respect to the standard basis.
        \end{enumerate}
        \end{exercise}
        
        \newpage
        \begin{exercise}
        Let $T:\mathbf{R}^3\to \mathbf{R}^3$ be the orthogonal projection to a $1$-dimensional linear subspace $L\subset \mathbf{R}^3$.
        \begin{enumerate}[label = (\arabic*)]
            \item List the eigenvalues of $T$.
            \item Write the characteristic polynomial $p_T(x)$ for $T$.
            \item Is $T$ diagonalizable? Briefly justify your answer.
        \end{enumerate}
        \end{exercise}
        
        \begin{exercise}
        Let $L$ be the line parameterized by $L(t)=(2t,-3t,t)$ for $t\in \mathbf{R}$, and let $T:\mathbf{R}^3\to \mathbf{R}^3$ be the linear transformation that is the orthogonal projection onto $L$.
        \begin{enumerate}[label = (\arabic*)]
            \item Describe $\operatorname{ker}(T)$ and $\operatorname{im}(T)$, either implicitly (using equations in $x,y,z$) or parametrically.
            \item List the eigenvalues of $T$ and their geometric multiplicities.
            \item Find a basis for each eigenspace of $T$.
            \item Let $A$ be the matrix for $T$ with respect to the standard basis. Find a diagonal matrix $B$ and an invertible matrix $S$ such that $B=S^{-1}AS$. (You do not have to compute $A$.)
        \end{enumerate}
        \end{exercise}
        
        \begin{exercise}
        Let $T:\mathbf{R}^4\to \mathbf{R}^4$ be orthogonal projection to the $2$-dimensional plane $P$ spanned by the vectors $\mathbf{v}=(2,0,1,0)$ and $\mathbf{w}=(-1,0,2,0)$.
        \begin{enumerate}[label = (\arabic*)]
            \item Find (with proof) all eigenvalues and eigenvectors, along with their geometric and algebraic multiplicities.
            \item Find the matrix representing $T$ with respect to the standard basis. Is this matrix diagonalizable? Why or why not?
        \end{enumerate}
        \end{exercise}
        
        \begin{exercise}
        Let $a, b \in \mathbf{R}$ and $T: \mathbf{R}^3 \to \mathbf{R}^3$ be the linear transformation that is the orthogonal projection onto the plane $z=ax+by$ (with respect to the usual Euclidean inner-product on $\mathbf{R}^3$).
        \begin{enumerate}[label = (\arabic*)]
            \item Find the eigenvalues of $T$ and bases for the corresponding eigenspaces.
            \item Is $T$ diagonalizable? Justify.
            \item What is the characteristic polynomial of $T$?
        \end{enumerate}
        \end{exercise}
        
        \begin{exercise}
        Let $T:\mathbf{R}^3\to \mathbf{R}^3$ be the orthogonal projection onto the plane $z=x+y$, with respect to the standard Euclidean inner product.
        \begin{enumerate}[label = (\arabic*)]
            \item Write the matrix representation of $T$ with respect to the standard basis.
            \item Is $T$ diagonalizable? Justify your answer.
        \end{enumerate}
        \end{exercise}
        \newpage
        
        \begin{exercise}
        Let $T:\mathbf{R}^3\to \mathbf{R}^3$ be the linear transformation that expands radially by a factor of three around the line parameterized by $L(t)=\begin{bmatrix} 2 \\ 2 \\ -1\end{bmatrix} t$, leaving the line itself fixed (viewed as a subspace).
        \begin{enumerate}[label = (\arabic*)]
            \item Find an eigenbasis for $T$ and provide the matrix representation of $T$ with respect to that basis.
            \item Provide the matrix representation of $T$ with respect to the standard basis.
        \end{enumerate}
        \end{exercise}
        
        \begin{exercise}
        Let $a,b\in \mathbf{R}$ and $T:\mathbf{R}^3\to \mathbf{R}^3$ be the linear transformation which is reflection across the plane $z=ax+by$.
        \begin{enumerate}[label = (\arabic*)]
            \item Find the eigenvalues of $T$ and for each find a basis for the corresponding eigenspace.
            \item Is $T$ diagonalizable? Justify.
            \item What is the characteristic polynomial of $T$?
            \item What is the minimal polynomial of $T$?
        \end{enumerate}
        \end{exercise}

%%%%%%%%%%%%%%%%%%%%%%%%%%%%%%%%%%%%%%%%%%%%%%%%%%%%%%%%%%%%%
%%%%%%%%%%%%%%%%%%%%%%%%%%%%%%%%%%%%%%%%%%%%%%%%%%%%%%%%%%%%%
%%%%%%%%%%%%%%%%%%%%%%%%%%%%%%%%%%%%%%%%%%%%%%%%%%%%%%%%%%%%%
%%%%%%%%%%%%%%%%%%%%%%%%%%%%%%%%%%%%%%%%%%%%%%%%%%%%%%%%%%%%%
%%%%%%%%%%%%%%%%%%%%%%%%%%%%%%%%%%%%%%%%%%%%%%%%%%%%%%%%%%%%%
%%%%%%%%%%%%%%%%%%%%%%%%%%%%%%%%%%%%%%%%%%%%%%%%%%%%%%%%%%%%%
%%%%%%%%%%%%%%%%%%%%%%%%%%%%%%%%%%%%%%%%%%%%%%%%%%%%%%%%%%%%%
%%%%%%%%%%%%%%%%%%%%%%%%%%%%%%%%%%%%%%%%%%%%%%%%%%%%%%%%%%%%%
%%%%%%%%%%%%%%%%%%%%%%%%%%%%%%%%%%%%%%%%%%%%%%%%%%%%%%%%%%%%%
%%%%%%%%%%%%%%%%%%%%%%%%%%%%%%%%%%%%%%%%%%%%%%%%%%%%%%%%%%%%%
%%%%%%%%%%%%%%%%%%%%%%%%%%%%%%%%%%%%%%%%%%%%%%%%%%%%%%%%%%%%%
%%%%%%%%%%%%%%%%%%%%%%%%%%%%%%%%%%%%%%%%%%%%%%%%%%%%%%%%%%%%%
%%%%%%%%%%%%%%%%%%%%%%%%%%%%%%%%%%%%%%%%%%%%%%%%%%%%%%%%%%%%%
%%%%%%%%%%%%%%%%%%%%%%%%%%%%%%%%%%%%%%%%%%%%%%%%%%%%%%%%%%%%%
%%%%%%%%%%%%%%%%%%%%%%%%%%%%%%%%%%%%%%%%%%%%%%%%%%%%%%%%%%%%%
%%%%%%%%%%%%%%%%%%%%%%%%%%%%%%%%%%%%%%%%%%%%%%%%%%%%%%%%%%%%%
%%%%%%%%%%%%%%%%%%%%%%%%%%%%%%%%%%%%%%%%%%%%%%%%%%%%%%%%%%%%%
%%%%%%%%%%%%%%%%%%%%%%%%%%%%%%%%%%%%%%%%%%%%%%%%%%%%%%%%%%%%%
\chapter*{Group Theory}
\section*{\S\ Pool Problems}
\markboth{Group Theory}{Pool Problems}
    \begin{exercise}
        Let $G$ be a group. Prove that $G$ is non-cyclic if and only if $G$ is the union of its proper subgroups.
        \end{exercise}
        
        \begin{exercise}
        Let $G$ be a group, and $G\times G$ the direct product. The set $D=\{(g,g)\mid g\in G\}$ is a subgroup of $G\times G$. Prove that if $D$ is normal in $G\times G$ then $G$ is abelian.
        \end{exercise}
        
        \begin{exercise}
        The dihedral group, $D_8$, is the group of eight rigid symmetries of a square. Prove that $D_8$ is not the internal direct product of two of its proper subgroups.
        \end{exercise}
        
        \begin{exercise}
        Let $G$ be a finite group and $H,K\mathrel{\unlhd}G$ be normal subgroups of relatively prime order. Prove that $G$ is isomorphic to a subgroup of $G/H\times G/K$.
        \end{exercise}
        
        \begin{exercise}
        Suppose $G$ is a group that contains normal subgroups $H,K\unlhd G$ with $H\cap K=\{e\}$ and $HK=G$. Prove that $G\cong H\times K$.
        \end{exercise}
        
        \begin{exercise}
        Let $G$ be the group of upper-triangular real matrices $\begin{bmatrix} a & b \\ 0 & d\end{bmatrix}$ with $a,d\neq 0$, under matrix multiplication. Let $S$ be the subset of $G$ defined by $d=1$. Show that $S$ is normal and that $G/S\cong \mathbf{R}^{\times}$, the multiplicative group of nonzero real numbers.
        \end{exercise}
        
        \begin{exercise}
        Let $G$ be a group and suppose $\operatorname{Aut}(G)$ is trivial.
        \begin{enumerate}[label = (\arabic*)]
            \item Show that $G$ is abelian.
            \item Show that for any abelian group $H$, the \textbf{inversion map} $\phi(h)=h^{-1}$ is an automorphism.
            \item Use parts (1) and (2) above to show that $g^2$ is the identity element for every $g\in G$.
        \end{enumerate}
        \end{exercise}
        
        \begin{exercise}
            \phantom{a}
        \begin{enumerate}[label = (\arabic*)]
            \item Suppose $N$ is a normal subgroup of a group $G$ and $\pi_N:G\to G/N$ is the usual projection homomorphism, defined by $\pi_N(g)=gN$. Prove that if $\phi:G\to H$ is any homomorphism with $N\leq \ker(\phi)$, then there exists a unique homomorphism $\psi:G/N\to H$ such that $\phi = \psi\circ \pi_N$. (You must explicitly define $\psi$, show it is well defined, show $\phi=\psi\circ\pi_N$, and show that $\psi$ is uniquely determined.)
            \item Prove the 
            \textbf{Third Isomorphism Theorem}: if $M, N\unlhd G$ with $N\le M$, then $(G/N)/(M/N)\cong G/M$.
        \end{enumerate}
        \end{exercise}
        
        \begin{exercise}
        Let $G$ be a group and $a\in G$ be an element. Let $n\in \mathbf{N}$ be the smallest positive number such that $a^n=e$, where $e$ is the identity element. Show that the set
        \begin{equation*}
        \begin{split}
        \{e,a,a^2,\ldots, a^{n-1}\}
        \end{split}
        \end{equation*}
        contains no repetitions.
        \end{exercise}
        
        \begin{exercise}
        Let $G$ be a finite abelian group of odd order. Let $\phi:G\to G$ be the function defined by $\phi(g)=g^2$ for all $g\in G$. Prove that $\phi$ is an automorphism.
        \end{exercise}
        
        \begin{exercise}
        Let $G$ be a group with exactly two conjugacy classes. Prove that $G$ is abelian, and describe all such groups (with proof).
        \end{exercise}
        
        \begin{exercise}
        Let $\mathbf{Z}_n$ denote the cyclic group of order $n$. Suppose $m\in \mathbf{N}$ is relatively prime to $n$. Define the function $\mu_m:\mathbf{Z}_n\to \mathbf{Z}_n$ by $m[a]_n=[ma]_n$.
        \begin{enumerate}[label = (\arabic*)]
            \item Prove that the map $\mu_m$ is a well-defined automorphism of $\mathbf{Z}_n$.
            \item Prove that any automorphism of $\mathbf{Z}_n$ has the form $\mu_m$ for some $m$.
        \end{enumerate}
        \end{exercise}
        
        \begin{exercise}
        For a group $G$ and an element $g\in G$, the \textbf{centralizer} of $g$ in $G$ is the subgroup
        \begin{equation*}
        \begin{split}
        C_G(g)=\{h\in G:hgh^{-1}=g\}.
        \end{split}
        \end{equation*}
        We say $g$ and $g'$ are \textbf{conjugate in $G$} if there exists an element $h\in G$ such that $g'=hgh^{-1}$.
        
        Suppose $S_n$ is a symmetric group with $n\ge 4$, and $\sigma$ is one of the $(n-2)$-cycles in $S_n$. (There are $\frac{n!}{2(n-2)}$ such cycles.)
        \begin{enumerate}[label = (\arabic*)]
            \item Prove that $[S_n:C_{S_n}(\sigma)]=[A_n:C_{A_n}(\sigma)]$.
            \item Determine whether all $(n-2)$-cycles are conjugate in $A_n$.
        \end{enumerate}
        \end{exercise}
        
        \begin{exercise}
        Let $G$ be a finite group and $n>1$ an integer such that $(ab)^n=a^n b^n$ for all $a,b\in G$. Let
        \begin{equation*}
        \begin{split}
        G_n=\{c\in G\mid c^n=e\},\qquad G^n=\{c^n\mid c\in G\}.
        \end{split}
        \end{equation*}
        You may take for granted that these are subgroups. Prove that both $G_n$ and $G^n$ are normal in $G$, and $|G^n|=[G:G_n]$.
        \end{exercise}
        
        \begin{exercise}
        Suppose $G$ is a group, $H\le G$ a subgroup, and $a,b\in G$. Prove that the following are equivalent:
        \begin{enumerate}[label = (\arabic*)]
            \item $aH=bH$
            \item $b\in aH$
            \item $b^{-1}a\in H$
        \end{enumerate}
        \end{exercise}
        
        \newpage
        \begin{exercise}
        Let $G$ be a group, and let $\operatorname{Aut}(G)$ denote the group of automorphisms of $G$. There is a homomorphism $\gamma:G\to \operatorname{Aut}(G)$ that takes $s\in G$ to the automorphism $\gamma_s$ defined by $\gamma_s(t)=sts^{-1}$.
        \begin{enumerate}[label = (\arabic*)]
            \item Prove rigorously, possibly with induction, that if $\gamma_s(t)=t^b$, then $\gamma_{s^n}(t)=t^{b^n}$.
            \item Suppose $s\in G$ has order 5, and $sts^{-1}=t^2$. Find the order of $t$. Justify your answer.
        \end{enumerate}
        \end{exercise}
        
        \begin{exercise}
        Let $G$ be an abelian group and $G_T$ be the set of elements of finite order in $G$.
        \begin{enumerate}[label = (\arabic*)]
            \item Prove that $G_T$ is a subgroup of $G$.
            \item Prove that every non-identity element of $G/G_T$ has infinite order.
            \item Characterize the elements of $G_T$ when $G=\mathbf{R}/\mathbf{Z}$, where $\mathbf{R}$ is the additive group of real numbers.
        \end{enumerate}
        \end{exercise}
        
        \begin{exercise}
        Suppose $G$ is a finite group of even order.
        \begin{enumerate}[label = (\arabic*)]
            \item Prove that an element in $G$ has order dividing 2 if and only if it is its own inverse.
            \item Prove that the number of elements in $G$ of order 2 is odd.
            \item Use (2) to show $G$ must contain a subgroup of order 2.
        \end{enumerate}
        \end{exercise}
        
        \begin{exercise}
        Let $N$ be a finite normal subgroup of $G$. Prove there is a normal subgroup $M$ of $G$ such that $[G:M]$ is finite and $nm=mn$ for all $n\in N$ and $m\in M$.
        
        \medskip
        \noindent ({Hint:} You may use the fact that the centralizer $C(h):=\{g\in G\mid ghg^{-1}=h\}$ is a subgroup of $G$.)
        \end{exercise}
        
        \begin{exercise}
        Show that every finite group with more than two elements has a nontrivial automorphism.
        \end{exercise}
        
        \begin{exercise}
        Suppose $G_1$ and $G_2$ are groups, with identity elements $e_1$ and $e_2$, respectively. Prove that if $\phi:G_1\to G_2$ is a homomorphism, then $\phi(e_1)=e_2$.
        \end{exercise}
        
        \begin{exercise}
        Suppose $A$ and $B$ are subgroups of a group $G$, and suppose $B$ is of finite index in $G$.
        \begin{enumerate}[topsep=0.1in,label = (\arabic*)]
            \item Show that the index of $A\cap B\le A$ is finite, and in fact $|A:A\cap B|\le |G:B|$. {Hint:} Find a set map $A/A\cap B\to G/B$.
            \item Prove that equality holds in (1) if and only if $G=AB$.
        \end{enumerate}
        \end{exercise}
        
        \begin{exercise}
        Let $G$ be a group. For each $a\in G$, let $\gamma_a$ denote the automorphism of $G$ defined by $\gamma_a(b)=aba^{-1}$ for all $b\in G$. The set $\operatorname{Inn}(G)=\{\gamma_a:a\in G\}$ is a subgroup of the automorphism group of $G$, called the subgroup of \textbf{inner automorphisms}. Prove that $\operatorname{Inn}(G)$ is isomorphic to $G/Z(G)$, where $Z(G)$ is the center of $G$.
        \end{exercise}
        
        \begin{exercise}
        Let $G$ be a group of order $2p$, where $p$ is an odd prime. Prove $G$ contains a nontrivial, proper normal subgroup.
        \end{exercise}
        
        \begin{exercise}
        Prove from the definition along that there are no nonabelian groups of order less than $5$.
        \end{exercise}
        
        \begin{exercise}
        Let $G$ be a group and $H,K\mathrel{\unlhd}G$ be normal subgroups with $H\cap K=\{e\}$. Show that each element in $H$ commutes with every element in $K$.
        \end{exercise}
        
        \begin{exercise}
        Let $G$ be a group and $N$ a normal subgroup of $G$. Let $aN$ denote the left coset defined by $a\in G$, and consider the binary operation
        \[
        G/N\times G/N\to G/N
        \]
        given by $(aN, bN)\mapsto abN$.
        \begin{enumerate}[label = (\arabic*)]
            \item Show the operation is well defined.
            \item Show the operation is well defined only if the subgroup $N$ is normal.
        \end{enumerate}
        \end{exercise}
        
        \begin{exercise}
        Let $G$ be a group, $H\le G$ a subgroup that is not normal. Prove there exist cosets $Ha$ and $Hb$ such that $HaHb\neq Hab$.
        \end{exercise}
        
        \begin{exercise}
        Let $H$ be a subgroup of a group $G$. The \textbf{normalizer} of $H$ in $G$ is the set $\mathbf{N}_G(H)=\{g\in G\,\mid\, gH=Hg\}$.
        \begin{enumerate}[label = (\arabic*)]
            \item Prove $\mathbf{N}_G(H)$ is a subgroup of $G$ containing $H$.
            \item Prove $\mathbf{N}_G(H)$ is the largest subgroup of $G$ in which $H$ is normal.
        \end{enumerate}
        \end{exercise}
        
        \begin{exercise}
        Let $G$ be a group and suppose $H\le G$. The \textbf{normalizer} of $H$ in $G$ is defined to be $N(H)=\{g\in G\,|\, gH=Hg\}$ and the \textbf{centralizer} of $H$ in $G$ is defined to be $C(H)=\{g\in G\,|\,  gh=hg\text{ for all }h\in H\}$.
        \begin{enumerate}[label = (\arabic*)]
            \item Prove that $N(H)$ is a subgroup of $G$.
            \item Prove that $C(H)$ is a normal subgroup of $N(H)$ and that $N(H)/C(H)$ is isomorphic to a subgroup of $\operatorname{Aut}(H)$.
        \end{enumerate}
        \end{exercise}

        \begin{exercise}
            Suppose $G$ is a cyclic group of finite order $n$, and $t\in G$ is a generator.
            \begin{enumerate}[label = (\arabic*)]
                \item Give a positive integer $d$ such that $t^{-1}=t^d$.
                \item Let $c$ be an integer and let $m=\gcd(n,c)$. Prove that the order of $t^c$ is $\frac{n}{m}$.
            \end{enumerate}
            \end{exercise}
        
        \begin{exercise}
        Let $G$ be a finite group. Prove \textit{from the definitions} that there exists a number $N$ such that $a^N=e$ for all $a\in G$.
        \end{exercise}
        
        \begin{exercise}
        Suppose $G$ is a group and $N\unlhd G$ is a finite normal subgroup. Prove that if $G/N$ contains an element of order $n$, then $G$ also contains an element of order $n$.
        \end{exercise}
        
        \begin{exercise}
        Suppose $\phi:G\to G'$ is a surjective homomorphism, $H\le G$ is a subgroup containing $\ker(\phi)$, and $H'=\phi(H)$. Prove $\phi^{-1}(H')=H$, where $\phi^{-1}(H')=\{g\in G\,\mid\, \phi(g)\in H'\}$. Make sure to state explicitly where each hypothesis is used.
        \end{exercise}
        
        \begin{exercise}
        Let $G$ be a group, and $H, K$ be subgroups of $G$. Let $HK=\{hk\,\mid \, h\in H, k\in K\}$ denote the set product. Prove that $HK$ is a group if and only if $HK=KH$.
        \end{exercise}
        
        \begin{exercise}
        Suppose $G$ is a nontrivial finite group and $H,K\mathrel{\unlhd}G$ are normal subgroups with $\gcd(|H|,|K|)=1$.
        \begin{enumerate}[label = (\arabic*)]
            \item Define a nontrivial group homomorphism $\phi:G\to G/H\times G/K$
            \item Prove $G$ is isomorphic to a subgroup of $G/H\times G/K$.
            \item Suppose $\gcd(m,n)=1$. Prove $\mathbf{Z}_{mn}\cong \mathbf{Z}_m\times \mathbf{Z}_n$.
        \end{enumerate}
        \end{exercise}
        
        \begin{exercise}
        Suppose $G$ is a group, $H$ and $K$ are normal subgroups of $G$, and $H\le K$.
        \begin{enumerate}[label = (\arabic*)]
            \item Define a group homomorphism from $K$ to $G/H$.
            \item Compute the kernel of the homomorphism in (1), and apply the First Isomorphism Theorem.
        \end{enumerate}
        \end{exercise}
        
        \begin{exercise}
        Let $G$ be a finite group and $\operatorname{Z}(G)$ denote its center.
        \begin{enumerate}[label = (\arabic*)]
            \item Prove that if $G/\operatorname{Z}(G)$ is cyclic, then $G$ is abelian.
            \item Prove that if $G$ is nonabelian, then $|\operatorname{Z}(G)|\le \frac{1}{4}|G|$.
        \end{enumerate}
        \end{exercise}
        
        \begin{exercise}
        Let $G$ be a group, $m\in \mathbf{N}$, and $g\in G$ be an element such that $g^m=e$. Prove that $\operatorname{o}(g)\mid m$, where $\operatorname{o}(g)$ is the order of $g$.
        \end{exercise}
        
        \begin{exercise}
            \phantom{a}
        \begin{enumerate}[label = (\arabic*)]
            \item Show that if $G$ is any group (not necessarily finite) and $H$ is a subgroup, then $G$ is a disjoint union of left cosets of $H$.
            \item State and prove Lagrange's Theorem for finite groups.
        \end{enumerate}
        \end{exercise}
        
        \begin{exercise}
        Let $G$ be a group and $H\le G$ a subgroup. For each coset $aH$ of $H$ in $G$, define the set
        \begin{equation*}
        \begin{split}
        G_{aH}=\{b\in G\,|\, baH=aH\}.
        \end{split}
        \end{equation*}
        \begin{enumerate}[label = (\arabic*)]
            \item Prove that $G_{aH}$ is a subgroup of $G$.
            \item Suppose that $H$ is normal in $G$. Prove that $G_{aH}=H$.
        \end{enumerate}
        \end{exercise}
        
        \begin{exercise}
        Let $G$ be a group of order $2n$ for some positive integer $n>1$.
        \begin{enumerate}[label = (\arabic*)]
            \item Prove there exists a subgroup $K$ of $G$ of order $2$.
            \item Suppose $K$ in (1) is a \textit{normal} subgroup. Prove that $K$ is contained in the center $\operatorname{Z}(G)$. (Recall $\operatorname{Z}(G)=\{a\in G\mid ab=ba\text{ for all }b\in G\}$.)
        \end{enumerate}
        \end{exercise}
        
        \begin{exercise}
        The additive group $\mathbf{Z}=(\mathbf{Z},+)$ of rational integers is a subgroup of the additive group $\mathbf{Q}=(\mathbf{Q},+)$. Show that $\mathbf{Z}$ has infinite index in $\mathbf{Q}$.
        \end{exercise}
%%%%%%%%%%%%%%%%%%%%%%%%%%%%%%%%%%%
%%%%%%%%%%%%%%%%%%%%%%%%%%%%%%%%%%%
%%%%%%%%%%%%%%%%%%%%%%%%%%%%%%%%%%%
%%%%%%%%%%%%%%%%%%%%%%%%%%%%%%%%%%%
%%%%%%%%%%%%%%%%%%%%%%%%%%%%%%%%%%%
%%%%%%%%%%%%%%%%%%%%%%%%%%%%%%%%%%%
%%%%%%%%%%%%%%%%%%%%%%%%%%%%%%%%%%%
%%%%%%%%%%%%%%%%%%%%%%%%%%%%%%%%%%%
%%%%%%%%%%%%%%%%%%%%%%%%%%%%%%%%%%%
%%%%%%%%%%%%%%%%%%%%%%%%%%%%%%%%%%%
%%%%%%%%%%%%%%%%%%%%%%%%%%%%%%%%%%%
%%%%%%%%%%%%%%%%%%%%%%%%%%%%%%%%%%%
%%%%%%%%%%%%%%%%%%%%%%%%%%%%%%%%%%%
%%%%%%%%%%%%%%%%%%%%%%%%%%%%%%%%%%%
%%%%%%%%%%%%%%%%%%%%%%%%%%%%%%%%%%%
%%%%%%%%%%%%%%%%%%%%%%%%%%%%%%%%%%%
%%%%%%%%%%%%%%%%%%%%%%%%%%%%%%%%%%%
%%%%%%%%%%%%%%%%%%%%%%%%%%%%%%%%%%%
\section*{\S\ Template Problems}

\markboth{Group Theory}{Template Problems}

        \begin{exercise}
            Let $G$ and $H$ be groups of order 10 and 15, respectively. Prove that if there is a nontrivial homomorphism $\phi:G\to H$, then $G$ is abelian.
            \end{exercise}
            
            \begin{exercise}
            Let $n$ be a number between $0$ and $10$. Compute $n^{111}\pmod{11}$, expressing your answer as a number between $0$ and $10$. Give as detailed a proof as you can, justifying every step, no matter how trivial you think it is.
            \end{exercise}
            
            \begin{exercise}
            Let $S_n$ denote the symmetric group on $n$ letters.
            \begin{enumerate}[label = (\arabic*)]
                \item Is the element $(1\,2\,3\,4)(2\,5\,3\,4\,6)(1\,5\,3\,2\,4\,7)\in S_7$ even or odd? Indicate your reasoning.
                \item Find the order of $(1\,3\,4)(2\,4\,3)(1\,3\,4)\in S_4$. Show all work.
                \item Write $(1\,5\,2\,3)(2\,1\,3\,4)(1\,5\,2\,3)^{-1}\in S_5$ in disjoint cycle form. Show all work.
            \end{enumerate}
            \end{exercise}
            
            \begin{exercise}
            Determine with proof the automorphism group $\operatorname{Aut}(V)$ of the Klein 4-group $V=\{e,a,b,ab\}$. To what familiar group is it isomorphic?
            \end{exercise}
            
            \begin{exercise}
            Determine the number of group homomorphisms $\phi$ between the given groups. Here $K_4$ denotes the Klein four-group (also known as $\mathbf{Z}/2\mathbf{Z}\times \mathbf{Z}/2\mathbf{Z}$) and $S_3$ denotes the symmetric group on three elements.
            \medskip
            \begin{enumerate}[label = (\arabic*)]
                \item $\phi:K_4\to \mathbf{Z}/2\mathbf{Z}$
                \item $\phi:\mathbf{Z}/2\mathbf{Z}\to K_4$
                \item $\phi:S_3\to K_4$
                \item $\phi:K_4\to S_3$
            \end{enumerate}
            \end{exercise}

            \begin{exercise}
                Without using Cauchy's Theorem or the Sylow theorems, prove that every group of order 21 contains an element of order three.
            \end{exercise}
            
            \begin{exercise}
            Explicitly list all group homomorphisms $f:\mathbf{Z}/6\mathbf{Z} \to \mathbf{Z}/12\mathbf{Z}$.
            \end{exercise}
            
            \begin{exercise}
            Let $C$ be a (possibly infinite) cyclic group, and let $\operatorname{Aut}(C)$ and $\operatorname{Inn}(C)$ be the groups of automorphisms and inner automorphisms, respectively. (Recall an automorphism $\gamma$ is \textbf{inner} if it is given by conjugation: $\gamma(b)=aba^{-1}$ for some $a\in C$.)
            \begin{enumerate}[label = (\arabic*)]
                \item Describe $\operatorname{Aut}(C)$ and $\operatorname{Inn}(C)$ in familiar terms, as groups you would study in a first algebra course. Prove your result. (\textit{Hint:} Where do generators go?)
                \item Write $\operatorname{Aut}(\mathbf{Z}_{12})$ down explicitly, giving its generic name and computing the order of every element. Show all work.
            \end{enumerate}
            \end{exercise}
            
            \newpage
            \begin{exercise}
            Let $A_5$ denote the alternating group on a $5$-element set $\{1,2,3,4,5\}$. The set of automorphisms of $A_5$ form a group, denoted $\operatorname{Aut}(A_5)$. The group of \textbf{conjugations} of $A_5$, denoted $\operatorname{Conj}(A_5)$, is the subgroup of $\operatorname{Aut}(A_5)$ consisting of automorphisms of the form $\gamma_s:=s(-)s^{-1}$ where $s\in A_5$. Explicitly, $\gamma_s(x)=sxs^{-1}$ for any $x\in A_5$.
            \begin{enumerate}[label = (\arabic*)]
                \item Prove that the function $\gamma:A_5\to \operatorname{Conj}(A_5)$, taking $s\in A_5$ to $\gamma_s$, is a surjective homomorphism.
                \item Prove that $A_5$ is isomorphic to $\operatorname{Conj}(A_5)$.
            \end{enumerate}
            \end{exercise}
            
            \begin{exercise}
            Suppose $H$ is a group of order 15. Prove there does not exist a nontrivial group homomorphism $\phi:D_5\to H$, where $D_5$ is the dihedral group with ten elements.
            \end{exercise}
            
            \begin{exercise}
            Let $S_7$ denote the symmetric group.
            \begin{enumerate}[label = (\arabic*)]
                \item Give an example of two non-conjugate elements of $S_7$ that have the same order.
                \item If $g\in S_7$ has maximal order, what is the order of $g$?
                \item Does the element $g$ that you found in part (2) lie in $A_7$? Fully justify your answer.
                \item Determine whether the set $\{h\in S_7 \mid |h|=|g|\}$ is a single conjugacy class in $S_7$, where $g$ is the element you found in part (2).
            \end{enumerate}
            \end{exercise}
            
            \begin{exercise}
            Let $G$ be the additive group $\mathbf{Z}_{2020}$ and let $H\subseteq G$ be the subset consisting of those elements with order dividing 20.
            \begin{enumerate}[label = (\arabic*)]
                \item Prove $H$ is a subgroup of $G$.
                \item Find an explicit generator for $H$ and determine its order.
            \end{enumerate}
            \end{exercise}
            
            \begin{exercise}
            Let $G$ denote the set of invertible $2\times 2$ matrices with values in a field. Prove $G$ is a group by defining a group law, identity element, and verifying the axioms. Credit is based on completeness.
            \end{exercise}
%%%%%%%%%%%%%%%%%%%%%%%%%%%%%%%%%%%
%%%%%%%%%%%%%%%%%%%%%%%%%%%%%%%%%%%
%%%%%%%%%%%%%%%%%%%%%%%%%%%%%%%%%%%
%%%%%%%%%%%%%%%%%%%%%%%%%%%%%%%%%%%
%%%%%%%%%%%%%%%%%%%%%%%%%%%%%%%%%%%
%%%%%%%%%%%%%%%%%%%%%%%%%%%%%%%%%%%
%%%%%%%%%%%%%%%%%%%%%%%%%%%%%%%%%%%
%%%%%%%%%%%%%%%%%%%%%%%%%%%%%%%%%%%
%%%%%%%%%%%%%%%%%%%%%%%%%%%%%%%%%%%
%%%%%%%%%%%%%%%%%%%%%%%%%%%%%%%%%%%
%%%%%%%%%%%%%%%%%%%%%%%%%%%%%%%%%%%
%%%%%%%%%%%%%%%%%%%%%%%%%%%%%%%%%%%
%%%%%%%%%%%%%%%%%%%%%%%%%%%%%%%%%%%
%%%%%%%%%%%%%%%%%%%%%%%%%%%%%%%%%%%
%%%%%%%%%%%%%%%%%%%%%%%%%%%%%%%%%%%
%%%%%%%%%%%%%%%%%%%%%%%%%%%%%%%%%%%
%%%%%%%%%%%%%%%%%%%%%%%%%%%%%%%%%%%
%%%%%%%%%%%%%%%%%%%%%%%%%%%%%%%%%%%
\chapter*{Ring Theory}
\section*{\S\ Pool Problems}
\markboth{Ring Theory}{Pool Problems}
\begin{exercise}
    Consider the additive group of integers $\mathbf{Z}$.
    \begin{enumerate}[label = (\arabic*)]
        \item Prove that every subgroup of $\mathbf{Z}$ is cyclic.
        \item Prove that every homomorphic image of $\mathbf{Z}$ is cyclic.
        \item Consider the \textit{ring} $\mathbf{Z}$. Exhibit a prime ideal of $\mathbf{Z}$ that is not maximal.
    \end{enumerate}
    \end{exercise}
    
    \begin{exercise}
    Let $R$ be an integral domain. Suppose $a$ and $b$ are non-associate irreducible elements in $R$, and the ideal $(a,b)$ generated by $a$ and $b$ is a proper ideal. Show that $R$ is not a principal ideal domain (PID).
    \end{exercise}
    
    \begin{exercise}
    Let $R$ be a commutative ring with identity. Suppose that for every $a \in R$ there is an integer $n \geq 2$ such that $a^n = a$. Show that every prime ideal of $R$ is maximal.
    \end{exercise}
    
    \begin{exercise}
    Let $R$ be a commutative ring with $1$, and $\sigma:R\to R$ be a ring automorphism.
    \begin{enumerate}[label = (\arabic*)]
        \item Show that $F=\{r\in R\mid \sigma(r)=r\}$ is a subring of $R$ with $1$.
        \item Show that if $\sigma^2$ is the identity map on $R$, then each element of $R$ is the root of a monic polynomial of degree 2 in $F[x]$, where $F$ is as in (1).
    \end{enumerate}
    \end{exercise}
    
    \begin{exercise}
    Suppose $R$ is a ring such that $r^2=r$ for every $r\in R$.
    \begin{enumerate}[label = (\arabic*)]
        \item Prove that $r=-r$ for every $r\in R$.
        \item Show that $R$ must be commutative. \textit{Hint: Consider $(a+b)^2$.}
    \end{enumerate}
    \end{exercise}
    
    \begin{exercise}
    Let $R$ be a commutative ring with $1$. The \textbf{characteristic} $\operatorname{char}(R)$ of $R$ is the unique integer $n\ge 0$ such that $\langle n\rangle \subset \mathbf{Z}$ is the kernel of the homomorphism $\theta:\mathbf{Z}\to R$ defined by
    \[
    \theta(m)=\begin{cases} \underbrace{1_R+\cdots +1_R}_{m}, & m\ge 0\\
    \underbrace{-1_R+\cdots + -1_R}_{|m|}, & m<0 \end{cases}.
    \]
    \begin{enumerate}[label = (\arabic*)]
        \item Prove that if $f:R\to S$ is a monomorphism of commutative rings with $1$, then $\operatorname{char}(R)=\operatorname{char}(S)$.
        \item Prove by giving an example that $\operatorname{char}(R)$ is not always preserved by ring homomorphism.
    \end{enumerate}
    \end{exercise}
    
    \newpage
    \begin{exercise}
    \phantom{a}
    \begin{enumerate}[label = (\arabic*)]
        \item Prove that for every commutative ring with unity $R$, there is a unique ring homomorphism $\phi_R:\mathbf{Z}\to R$, and that $\ker(\phi_R)=\langle d_R\rangle$ for a unique nonnegative integer $d_R$. The number $d_R$ is called the \textbf{characteristic} of $R$, denoted $\operatorname{char}(R)$.
        \item Suppose $F_1$ and $F_2$ are fields for which there exists a ring homomorphism $f:F_1\to F_2$. Prove that $\operatorname{char}(F_1)=\operatorname{char}(F_2)$.
    \end{enumerate}
    \end{exercise}
    
    \begin{exercise}
    Let $A$ be a commutative ring with $1$. The \textbf{dimension} of $A$ is the maximal length $d$ of a chain of prime ideals $\mathfrak{p}_0\subsetneq \mathfrak{p}_1\subsetneq \cdots \subsetneq \mathfrak{p}_d$. Prove that if $A$ is a PID, the dimension of $A$ is at most 1.
    \end{exercise}
    
    \begin{exercise}
    Prove that every Euclidean domain is a principal ideal domain.
    \end{exercise}
    
    \begin{exercise}
    Let $F$ be a field and let $\alpha$ generate a field extension of $F$ of degree 5. Prove that $\alpha^2$ generates the same extension.
    \end{exercise}
    
    \begin{exercise}
    Let $R$ be a commutative ring with $1$. Use theorems in ring theory to prove:
    \begin{enumerate}[label = (\arabic*)]
        \item $\langle x\rangle$ is a prime ideal in $R[x]$ if and only if $R$ is an integral domain.
        \item $\langle x\rangle$ is a maximal ideal in $R[x]$ if and only if $R$ is a field.
    \end{enumerate}
    \end{exercise}
    
    \begin{exercise}
    Let $F$ be a field and $F[x]$ the polynomial ring, which is a principle ideal domain. Let $R=\{f\in F[x]: f'\in (x)\}$, where $(x)\subset F[x]$ is the ideal generated by $x$, and $f'$ is the (formal) derivative of the polynomial $f$. It is a fact that $R$ is a subring of $F[x]$.
    \begin{enumerate}[label = (\arabic*)]
        \item Prove that $x^2$ and $x^3$ are irreducible elements of $R$.
        \item Let $(x^2,x^3)$ be the ideal in $R$ generated by $x^2$ and $x^3$. Prove it is a proper ideal of $R$.
        \item Prove that $(x^2,x^3)$ is not a \textit{principal} ideal of $R$.
    \end{enumerate}
    \end{exercise}
    
    \begin{exercise}
    Let $R$ be a commutative ring with $1$ and suppose $e \in R$ is \textbf{idempotent}; i.e., satisfies $e^2 = e$.
    \begin{enumerate}[label = (\arabic*)]
        \item Prove that $1-e$ is idempotent.
        \item Suppose $e \neq 0,1$. Show that $Re$ and $R(1-e)$ are proper ideals of $R$.
        \item Prove there is an isomorphism $R \cong Re \times R(1-e)$.
    \end{enumerate}
    \end{exercise}
    
    \begin{exercise}
    An element $r$ of a ring $R$ is said to be \textbf{idempotent} if $r^2 = r$. Suppose that $R$ is a commutative ring with unity containing an idempotent element $e$.
    \begin{enumerate}[label = (\arabic*)]
        \item Prove that $1-e$ is also idempotent.
        \item Prove that $eR$ and $(1-e)R$ are both ideals in $R$ and that $R\cong eR\times (1-e)R$. 
        \item Prove that if $R$ has a unique maximal ideal, then the only idempotent elements in $R$ are $0$ and $1$.
    \end{enumerate}
    \end{exercise}
    
    \begin{exercise}
    Prove that if $\phi:R\to S$ is a surjective ring homomorphism between commutative rings with unity, then $\phi(1_R)=1_S$.
    \end{exercise}
    
    \begin{exercise}
    Suppose $R$ is a PID (principal ideal domain). Prove that an ideal $I\subset R$ is maximal if and only if $I=\langle p\rangle$ for a prime $p\in R$. (By definition, an element $p$ is \textbf{prime} if whenever $p \mid ab$ then $p \mid a$ or $p \mid b$. If you use the fact that prime implies irreducible, you have to prove it.)
    \end{exercise}
    
    \begin{exercise}
    Let $R$ be a commutative ring.
    \begin{enumerate}[label = (\arabic*)]
        \item Prove that the set $N$ of all nilpotent elements of $R$ is an ideal.
        \item Prove that $R/N$ is a ring with no nonzero nilpotent elements.
        \item Show that $N$ is contained in every prime ideal of $R$.
    \end{enumerate}
    \end{exercise}
    
    \begin{exercise}
    Let $R$ be a commutative ring with $1$. We say an element $n \in R$ is \textbf{nilpotent} if there exists a number $k \in \mathbf{N}$ such that $n^k = 0$.
    \begin{enumerate}[label = (\arabic*)]
        \item Show that if $n$ is nilpotent, then $1-n$ is a unit.
        \item Give an example of a commutative ring with $1$ that has no nonzero nilpotent elements, but is not an integral domain.
    \end{enumerate}
    \end{exercise}
    
    \begin{exercise}
    Let $D$ be a principal ideal domain. Prove that every proper nonzero prime ideal is maximal.
    \end{exercise}
    
    \begin{exercise}
    Let $I\subseteq \mathbf{Z}[x]$ denote the set of all polynomials with even constant term.
    \begin{enumerate}[label = (\arabic*)]
        \item Prove that $I$ is an ideal.
        \item Prove that $I$ is not a principal ideal.
    \end{enumerate}
    \end{exercise}
    
    \begin{exercise}
    Let $R$ be a commutative ring with unity.
    \begin{enumerate}[label = (\arabic*)]
        \item Define what it means for an element in $R$ to be \textbf{prime}, and also what it means for an element to be \textbf{irreducible}.
        \item Prove that if $R$ is an integral domain, then every prime element is irreducible.
    \end{enumerate}
    \end{exercise}
    
    \begin{exercise}
    Let $R$ be a commutative ring with unity, let $I\subseteq R$ be an ideal, and let $\pi:R\to R/I$ be the natural projection homomorphism.
    \begin{enumerate}[label = (\arabic*)]
        \item Show that if $\wp$ is a prime ideal of $R/I$, then $\pi^{-1}(\wp)$ is a prime ideal of $R$.
        \item Show that the assignment $\wp \mapsto \pi^{-1}(\wp)$ is injective on the set of prime ideals of $R/I$.
    \end{enumerate}
    \end{exercise}
    
    \newpage
    \begin{exercise}
    Let $A$ be a commutative ring with unit. We call $A$ \textbf{Boolean} if $a^2=a$ for every $a\in A$. Prove that in a Boolean ring $A$ each of the following holds:
    \begin{enumerate}[label = (\arabic*)]
        \item $2a=0$ for every $a\in A$.
        \item If $I$ is a prime ideal then $A/I$ has two elements (and in particular $I$ is maximal).
        \item If $I=(a,b)$ is the ideal generated by $a$ and $b$ then $I$ can be generated by the single element $a + b + ab$. Conclude that every finitely generated ideal is principal.
    \end{enumerate}
    \end{exercise}
    
    \begin{exercise}
    Let $R$ be a commutative ring. For each nonempty subset $X \subseteq \mathbf{R}$, the \textbf{annihilator} of $X$ is the set $\operatorname{ann}(X) = \{a \in R \mid ax = 0 \text{ for all }x \in X\}$.
    \begin{enumerate}[label = (\arabic*)]
        \item Prove that $\operatorname{ann}(X)$ is an ideal of $R$.
        \item Prove that $X\subseteq \operatorname{ann}(\operatorname{ann}(X))$.
    \end{enumerate}
    \end{exercise}
    
    \begin{exercise}
    Suppose $\phi:R\to S$ is a ring homomorphism, and $S$ has no (nonzero)zero-divisors. Prove from the definitions that $\ker(\phi)$ is a prime ideal.
    \end{exercise}
    
    \begin{exercise}
    Let $R$ be a commutative ring with $1$, and $N$ the ideal 
        \begin{equation*}
        \begin{split}
            N=\{a\in R\mid a^n=0 \text{ for some }n\}.
        \end{split}
        \end{equation*}
    Let $[b]$ be the image of $b\in R$ in $R/N$. Prove that if $[a] \in R/N$ and $[a]^m = 0$ then $[a] = [0]$.
    \end{exercise}
    
    \begin{exercise}
    Let $R$ be a commutative ring. The \textbf{nilradical} of $R$ is defined to be $N=\{r\in R\mid r^n=0 \text{ for some } n\in \mathbf{N}\}$.
    \begin{enumerate}[label = (\arabic*)]
        \item Prove that $N$ is an ideal of $R$.
        \item Prove that $N$ is contained in the intersection of all prime ideals of $R$.
    \end{enumerate}
    \end{exercise}

%%%%%%%%%%%%%%%%%%%%%%%%%%%%%%%%%%%
%%%%%%%%%%%%%%%%%%%%%%%%%%%%%%%%%%%
%%%%%%%%%%%%%%%%%%%%%%%%%%%%%%%%%%%
%%%%%%%%%%%%%%%%%%%%%%%%%%%%%%%%%%%
%%%%%%%%%%%%%%%%%%%%%%%%%%%%%%%%%%%
%%%%%%%%%%%%%%%%%%%%%%%%%%%%%%%%%%%
%%%%%%%%%%%%%%%%%%%%%%%%%%%%%%%%%%%
%%%%%%%%%%%%%%%%%%%%%%%%%%%%%%%%%%%
%%%%%%%%%%%%%%%%%%%%%%%%%%%%%%%%%%%
%%%%%%%%%%%%%%%%%%%%%%%%%%%%%%%%%%%
%%%%%%%%%%%%%%%%%%%%%%%%%%%%%%%%%%%
%%%%%%%%%%%%%%%%%%%%%%%%%%%%%%%%%%%
\section*{\S\ Template Problems}
\markboth{Ring Theory}{Template Problems}
\begin{exercise}
    \phantom{a}
    \begin{enumerate}[label = (\arabic*)]
        \item Suppose $I$ and $J$ are ideals in a commutative ring $R$ such that $R=I+J$. Prove that the map $f:R\to R/I\times R/J$ given by $f(x)=(x+I,x+J)$ induces the isomorphism
        \[
            R/IJ \cong R/I \times R/J.
        \]
        \item Prove that $\left(\mathbf{Z}/3\mathbf{Z}\right)[x]/(x^3-x^2-1)\cong \left(\mathbf{Z}/3\mathbf{Z}\right)[x]/(x^3+x+1)$. (\textit{Hint: Use part (1).})
    \end{enumerate}
    \end{exercise}
    
    \begin{exercise}
    Let $\mathcal{C}([0,1])$ be the (commutative) ring of continuous, real-valued functions on the unit interval, and let
    \[
        M=\{f\in \mathcal{C}([0,1]) \mid f(1/2)=0\}.
    \]
    Prove that $M$ is a maximal ideal.
    \end{exercise}
    
    \begin{exercise}
    Let $z \in \mathbf{C}$ be a complex number and let $\epsilon_z:\mathbf{R}[x]\to \mathbf{C}$ be the evaluation homomorphism given by $\epsilon_z(p)=p(z)$ for each $p \in \mathbf{R}[x]$.
    \begin{enumerate}[label = (\arabic*)]
        \item Show that $\ker(\epsilon_z)$ is a prime ideal.
        \item Compute $\ker(\epsilon_{1+i})$, $\operatorname{im}(\epsilon_{1+i})$ and then state the conclusion of the First Isomorphism Theorem applied to the homomorphism $\epsilon_{1+i}$.
    \end{enumerate}
    \end{exercise}
    
    \begin{exercise}
    Let $\mathbf{Z}[\sqrt{2}]=\{a+b\sqrt{2} \mid a,b\in \mathbf{Z}\}$, and let $R\subset \operatorname{M}_2(\mathbf{Z})$ be the ring of all $2\times 2$ matrices of the form $\begin{bmatrix} a & b \\ 2b & a\end{bmatrix}$. Prove $\mathbf{Z}[\sqrt{2}]$ is isomorphic to $R$. \textit{Hint:} Start by showing the map $a+b\sqrt{2} \mapsto \begin{bmatrix} a & b \\ 2b & a\end{bmatrix}$ is a ring homomorphism.
    \end{exercise}
    
    \begin{exercise}
    Let $f(x)=x^3+x+1\in \mathbf{Z}_5[x]$.
    \begin{enumerate}[label = (\arabic*)]
        \item Prove that $f(x)$ is irreducible.
        \item Prove that $\langle f(x)\rangle$ is a maximal ideal.
        \item What is the cardinality of $\mathbf{Z}_5[x]/\langle f(x)\rangle$? Justify.
    \end{enumerate}
    \end{exercise}
    
    \begin{exercise}
    \phantom{a}
    \begin{enumerate}[label = (\arabic*)]
        \item Write down an irreducible cubic polynomial over $\mathbf{F}_2$.
        \item Construct a field with exactly 8 elements and write down its multiplication table.
    \end{enumerate}
    \end{exercise}
    
    \begin{exercise}
    Let $\varepsilon:\mathbf{R}[x]\to \mathbf{C}$ be the ring homomorphism that is evaluation at $i$, so $\varepsilon(f) = f(i)$ (here $i$ denotes the complex numbers sometimes denoted $\sqrt{-1}$).
    \begin{enumerate}[label = (\arabic*)]
        \item Prove that $\ker(\varepsilon)=(x^2+1)\subseteq \mathbf{R}[x]$.
        \item Prove that $(x^2+1)$ is a maximal ideal in $\mathbf{R}[x]$.
    \end{enumerate}
    \end{exercise}
    
    \begin{exercise}
    Let $\mathbf{Z}[i]=\{a+bi \mid a,b\in \mathbf{Z},\, i^2=-1\}$, a subring of $\mathbb{C}$. Prove there is no ring homomorphism $\mathbf{Z}[i]\to \mathbf{Z}_{19}$, but there is a ring homomorphism $\mathbf{Z}[i]\to \mathbf{Z}_{13}$. Note a ring homomorphism of commutative rings with $1$ must send $1$ to $1$ (\textit{Hint:} The group of units in $\mathbf{Z}_{19}$ is the cyclic group $U(19)$ of order 18, and the group of units in $\mathbf{Z}_{13}$ is the cyclic group $U(13)$ of order 12).
    \end{exercise}
    
    \begin{exercise}
    Let $\mathbf{Z}_n$ be the ring of integers $\pmod{n}$. There is a ring homomorphism:
        \begin{equation*}
        \begin{split}
            \mathbf{Z}_{28} \rightarrow \mathbf{Z}_5 \times \mathbf{Z}_7, \hspace{7pt} [m]_{28} \mapsto ([m]_4 , [m]_7).
        \end{split}
        \end{equation*}
    This is an isomorphism by the Chinese Remainder Theorem. Let $\mathbf{Z}_n^\times$ be the group of units of $\mathbf{Z}_n$. Prove that $\mathbf{Z}_{28}^\times$ is isomorphism to $\mathbf{Z}_4^\times \times \mathbf{Z}_7^\times$.
    \end{exercise}
    
    \begin{exercise}
    Let $k \subset K$ be fields, and let $k[X]$ be the polynomial ring in one variable with coefficients in $k$. The \textbf{evaluation at $z \in K$} is a ring homomorphism $\varepsilon:k[X]\to K$ defined by $\varepsilon(f(X))=f(z)$. Prove that if $\varepsilon$ is not injective, then $\varepsilon(k[X])$ is a field.
    \end{exercise}
    
    \begin{exercise}
    Let $i$ be the imaginary number, let $\mathbf{Z}[i] = \{a = bi \mid a,b \in \mathbf{Z}\}$, a principal ideal domain, and let $\mathbf{Z}_2$ be the finite ring of integers modulo 2.
    \begin{enumerate}[label = (\arabic*)]
        \item Define a ring homomorphism $\mathbf{Z}[i]\to \mathbf{Z}_2$. You must prove it is a ring homomorphism.
        \item Find, with proof, a generator for the kernel of your ring homomorphism.
    \end{enumerate}
    \end{exercise}
    
    \begin{exercise}
    Let $i \in \mathbf{C}$ be the usual root of unity, with $i^2 = -1$, and let $Z[i] = \{a+bi \mid a,b \in \mathbf{Z}\}$ be the ring of Gaussian integers.
    \begin{enumerate}[label = (\arabic*)]
        \item Prove that there exists a (nonzero) ring homomorphism $\mathbf{Z}[i] \rightarrow \mathbf{Z}_5$.
        \item Compute the kernel of your homomorphism explicity, and state the conclusion given by the First Isomorphism Theorem.
    \end{enumerate}
    \end{exercise}
    
    \begin{exercise}
    Let $\mathbf{Z}_2 = \{0,1\}$ be the field of two elements. The quotient ring $\mathbf{Z}_2[x]/\langle x^3+x+1\rangle$ is a field of cardinality 8, containing $\mathbf{Z}_2$. Let $\pi:\mathbf{Z}_2[x] \rightarrow \mathbf{Z}_2[x]/\langle x^3+x+1\rangle$ be the natural projection.
    \begin{enumerate}[label = (\arabic*)]
        \item Write down a set of eight distinct coset representatives for the elements of this field.
        \item Determine the multiplicative inverse of $\pi(x)$ in terms of your coset representatives.
    \end{enumerate}
    \end{exercise}
    
    \begin{exercise}
    Let $\mathbf{F}_9$ denote the field of nine elements.
    \begin{enumerate}[label = (\arabic*)]
        \item Show that each nonzero $a\in \mathbf{F}_9$ is a root of $X^3-1=(X-1)(X^2+1)(X^4+1)\in \mathbf{F}_3[X]$.
        \item Use the Pigeonhole Principle to prove that $\mathbf{F}_9$ has an element of multiplicative order 8. (Include a proof that the Pigeonhole Principle applies.)
    \end{enumerate}
    \end{exercise}
    
    \begin{exercise}
    Let $\mathbf{Z}[X]$ be the ring of polynomials with integer coefficients, and let $K\subset \mathbf{Z}[X]$ be the kernel of the ``evaluation at 1'' homomorphism:
        \begin{equation*}
        \begin{split}
            \epsilon_1:\mathbf{Z}[X] \rightarrow \mathbf{Z}_3, \hspace{7pt} f(X) \mapsto [f(1)]_3.
        \end{split}
        \end{equation*}
        \begin{enumerate}[label = (\arabic*),itemsep=1pt,topsep=3pt]
            \item Characterize $K$ as a set.
            \item Determine whether $K$ is a maximal ideal. Fully justify your conclusion.
            \item Determine whether $K$ is a principal ideal. Justify by either exhibiting a generator or proving that there isn't one.
        \end{enumerate}
    \end{exercise}

\end{document}

